% The data set as delivered (1,801-patient table, Kit 2026-10-28) and the baseline CT volumes, input
% to the in-house tracer. Only the data set is described here; the restriction to the study cohort
% is in methodology/02_population_outcome.tex (sec:method-population).
% Follow-up interval: section 1 of the cgps_stats.py report on the 1,801-patient table.
\section{Study Population and CT Volumes}
\label{sec:data-cohort}

The Copenhagen General Population Study, referred to in this thesis as Herlev--{\O}sterbro, is a prospective cohort study of adults from the Greater Copenhagen area, randomly selected from the Danish Civil Registration System.
Within this cohort, a subset of the participants underwent CCTA, and those scanned again at a follow-up visit form the basis of this thesis.
The data set comprises 1,801 participants, scanned a median of [..] years apart (interquartile range [..]--[..] years).
The study cohort drawn from it is defined in Section~\ref{sec:method-population}.

At each visit, a non-contrast, electrocardiogram (ECG)-gated scan was acquired for coronary artery calcium scoring, followed by contrast injection and an ECG-gated CCTA.
Only the baseline CCTA volumes were analysed in this thesis, serving as input to the in-house tracer (Section~\ref{sec:method-tracer}) and providing the centerlines and lumen from which the geometric features were computed.
The follow-up scans contributed through their visual analysis (Section~\ref{sec:data-plaque-location}).
\clarify{ number of baseline volumes received; scanner, tube voltage, contrast and reconstruction (voxel size, kernel).}
